\itemhead{4.1 例题 · 例题2（教材p.96）}{XJ-B1-C04-EX-P096-Y0533}
\label{XJ-B1-C04-EX-P096-Y0533}
{\footnotesize 湘教版教材p.96 / PDF第103页 / 4.1 例题 / 例题2}
\field{原题}{求下列各式的值：
\subq{1}$16^{3/4}$；
\subq{2}$25^{-1/2}$；
\subq{3}$\left(\dfrac13\right)^{-3}$；
\subq{4}$\left(\dfrac{125}{64}\right)^{-2/3}$。}
\field{参考答案}{依次为 $8,1/5,27,16/25$。}
{\small \field{关联知识}{P02 有理数、实数指数幂与运算（教材帮p.144、145 / PDF153、154）}
\field{方法与易错点}{统一为指数幂，分开处理系数与各底数的指数；最终复查定义条件及数值精度。 易错：负指数误作负数；幂的幂指数相加；约分丢禁点；过早四舍五入。}
\field{答案依据}{例题原解位置：湘教版教材p.96及题库原图；独立编写可编辑参考解答。答案册此部分主要收练习、习题，不伪造例题答案册位置。}}
